\documentclass[11pt]{article}

\usepackage[a4paper,margin=1in]{geometry}
\usepackage{amsmath,amssymb,bm}
\usepackage{booktabs}
\usepackage{hyperref}

\hypersetup{
    colorlinks=true,
    linkcolor=blue,
    citecolor=blue,
    urlcolor=blue
}

\newcommand{\eps}{\varepsilon}
\newcommand{\mean}[1]{\langle #1 \rangle}
\newcommand{\master}{geqoe\_averaged\_j2\_zonal\_note}

\title{Drag Extension for the GEqOE Averaged Theory\\
\large Secular Drag as an Add-On Perturbation to the Zonal Baseline}
\author{Jose Manuel Montilla Garc\'ia}
\date{\today}

\begin{document}
\maketitle

\begin{abstract}
This companion note describes how atmospheric drag can be incorporated into
the GEqOE averaged zonal theory (master note, \texttt{\master{}}) as an
add-on secular perturbation. The key insight is that drag primarily affects
the generalized energy $\nu$ (secular decay) and the eccentricity magnitude
$g$ (secular circularization), both of which are exactly constant in the
conservative zonal theory. The drag-averaged secular rates can be added to
the mean ODE without rederiving the zonal short-period map, because the
short-period corrections are parameterized by the instantaneous mean state
$(\bar\nu, \bar g, \bar Q)$ and remain valid as these quantities evolve
slowly under drag. The main approximation is that drag short-period effects
(once-per-orbit drag modulation) are neglected, which is justified when the
drag perturbation is much smaller than the zonal perturbation in its
short-period content.
\end{abstract}

\tableofcontents

\section{Motivation}

The master note develops the averaged GEqOE theory for conservative,
autonomous, axisymmetric perturbations (Earth zonal harmonics). In this
setting, the generalized energy $\nu$ is an exact first integral
($\dot\nu = 0$), and the eccentricity magnitude $g$ is either constant
(pure $J_2$) or slowly varying through the argument of pericenter (general
zonals).

Atmospheric drag breaks both of these: it dissipates orbital energy
($\dot\nu \ne 0$) and circularizes the orbit ($\dot g < 0$ secularly).
However, drag is typically much weaker than the zonal perturbation in its
\emph{short-period} content---the once-per-orbit modulation of drag is small
compared to the zonal short-period corrections. This suggests a natural
separation:
\begin{itemize}
\item Use the zonal short-period map (unchanged) for the mean-to-osculating
conversion.
\item Add drag-averaged secular rates to the mean ODE for the long-term
evolution.
\end{itemize}

This is the standard approach in semi-analytical theories (e.g., DSST), and
it is well justified when the drag perturbation magnitude is much smaller
than the zonal perturbation magnitude, or when the drag short-period effects
are not needed.


\section{Drag in the GEqOE Framework}

\subsection{Non-conservative GEqOE equations}

For a general non-conservative acceleration
$\bm P = P_r \bm e_r + P_t \bm e_t + P_h \bm e_h$
in the radial-transverse-normal (RTN) frame, the GEqOE equations of motion
acquire additional terms beyond the conservative zonal dynamics. The most
important new feature is that $\nu$ is no longer constant:
\begin{equation}
\dot\nu = -\frac{2\nu^{4/3}}{\mu^{2/3}}
\left(P_r \frac{\dot r}{\sqrt{\mu/a}} + P_t \frac{h}{r\sqrt{\mu/a}}\right),
\label{eq:nudot_general}
\end{equation}
where $h$ is the angular momentum magnitude, $\dot r$ is the radial
velocity, and $a = (\mu/\nu^2)^{1/3}$.

The full non-conservative GEqOE system is given in the master note's
parent reference (Ba\`u et al., 2021). For the present purpose, only the
orbit-averaged effect of drag is needed.

\subsection{Drag acceleration model}

The standard atmospheric drag acceleration is
\begin{equation}
\bm P_{\mathrm{drag}}
= -\frac{1}{2}\,\frac{C_D A}{m}\,\rho(r)\,v_{\mathrm{rel}}\,
\bm v_{\mathrm{rel}},
\label{eq:drag_accel}
\end{equation}
where $C_D$ is the drag coefficient, $A/m$ is the area-to-mass ratio,
$\rho(r)$ is the atmospheric density (assumed spherically symmetric for the
present treatment), and $\bm v_{\mathrm{rel}}$ is the velocity relative to
the co-rotating atmosphere.

For simplicity, the initial treatment neglects atmospheric co-rotation
(setting $\bm v_{\mathrm{rel}} \approx \bm v$) and assumes a spherically
symmetric exponential density model:
\begin{equation}
\rho(r) = \rho_0 \exp\!\left(-\frac{r - r_0}{H}\right),
\label{eq:density_exp}
\end{equation}
where $\rho_0$ is the reference density at altitude $r_0$ and $H$ is the
scale height.


\section{Orbit-Averaged Drag Rates}

\subsection{Averaging procedure}

The orbit-averaged drag effect on any slow variable $z$ is
\begin{equation}
\mean{\dot z_{\mathrm{drag}}}_K
= \frac{1}{2\pi a}\int_0^{2\pi}
r(K)\,\dot z_{\mathrm{drag}}(K)\,dK,
\label{eq:drag_avg}
\end{equation}
using the same averaging operator as the zonal theory (master note,
Section~4).

\subsection{Averaged energy dissipation}

The most important drag effect is the secular decay of $\nu$ (equivalently,
the secular decay of $a$). Using the vis-viva relation and the
orbit-averaged power dissipation:
\begin{equation}
\mean{\dot\nu_{\mathrm{drag}}}
= \frac{2\nu}{3a}\,\mean{\dot a_{\mathrm{drag}}},
\label{eq:nudot_drag_avg}
\end{equation}
where
\begin{equation}
\mean{\dot a_{\mathrm{drag}}}
= -\frac{C_D A}{m}\,\frac{1}{2\pi a}
\int_0^{2\pi} r\,\rho(r)\,v^2\,dK.
\label{eq:adot_drag_avg}
\end{equation}

For a Keplerian ellipse with semi-major axis $a$ and eccentricity $e = g$
(to first order), the orbital velocity satisfies
$v^2 = \mu(2/r - 1/a)$ and the radius is
$r = a(1 - g\cos G)$ where $G = K - \Psi$.

The integral \eqref{eq:adot_drag_avg} with the exponential density
\eqref{eq:density_exp} does not admit a closed-form elementary evaluation
because of the $\rho(r) = \rho_0 \exp(-(r-r_0)/H)$ factor, which through
$r = a(1 - g\cos G)$ becomes $\exp(ag\cos G / H)$. The standard treatment
expands this exponential in modified Bessel functions:
\begin{equation}
\exp\!\left(\frac{ag\cos G}{H}\right)
= I_0\!\left(\frac{ag}{H}\right)
+ 2\sum_{k=1}^{\infty}
I_k\!\left(\frac{ag}{H}\right)\cos(kG),
\label{eq:bessel_expansion}
\end{equation}
where $I_k$ is the modified Bessel function of the first kind. After
orbit averaging, only the $k = 0$ term survives in the secular part, and the
$k = 1$ term contributes to the eccentricity evolution.

\subsection{King-Hele formulas}

The classical orbit-averaged drag rates, originally due to King-Hele, are
(for the simplified case of no atmospheric rotation):
\begin{align}
\mean{\dot a_{\mathrm{drag}}}
&= -\frac{C_D A}{m}\,\rho_p\,a\,n\,a\,
\exp(-c_H)\,
\left[
I_0(c_H) + 2e\,\frac{I_1(c_H)}{1}
\right]
\cdot \ldots
\label{eq:adot_KH}
\end{align}
where $\rho_p = \rho(r_p)$ is the density at periapsis,
$c_H = ae/H$ is the ratio of the apsidal altitude variation to the scale
height, and $n = \sqrt{\mu/a^3}$ is the mean motion.

The exact King-Hele formulas are somewhat involved and depend on the
approximation order retained in the Bessel expansion. For the present design
note, the important structural point is that the averaged drag rates are
\emph{known functions} of $(a, e, i)$ and the ballistic coefficient
$C_D A/m$, with the atmospheric model entering through $\rho_p$ and $H$.

\subsection{Drag-averaged rates in GEqOE variables}

Translating the King-Hele formulas to GEqOE variables:
\begin{align}
\mean{\dot\nu_{\mathrm{drag}}}
&= \frac{2\nu}{3a}\,\mean{\dot a_{\mathrm{drag}}}
= F_\nu(\bar\nu, \bar g; C_D A/m, \rho_0, H),
\label{eq:nudot_geqoe}\\
\mean{\dot g_{\mathrm{drag}}}
&= F_g(\bar\nu, \bar g; C_D A/m, \rho_0, H),
\label{eq:gdot_geqoe}\\
\mean{\dot Q_{\mathrm{drag}}} &= 0,
\label{eq:Qdot_drag}\\
\mean{\dot\Psi_{\mathrm{drag}}} &= 0,
\label{eq:Psidot_drag}\\
\mean{\dot\Omega_{\mathrm{drag}}} &= 0.
\label{eq:Omegadot_drag}
\end{align}

The key structural facts:
\begin{enumerate}
\item \textbf{Drag affects only $\nu$ and $g$} (to leading order). The
inclination-related variables $Q$, $\Psi$, $\Omega$ are unaffected by drag
(in the no-atmospheric-rotation approximation), because drag acts in the
velocity direction and does not exert an out-of-plane force.

\item \textbf{Both $\nu$ and $g$ decrease secularly}: $\nu$ decreases
(energy loss $\Rightarrow$ smaller $a$), and $g$ decreases (circularization).

\item \textbf{The drag rates depend on $\nu$ and $g$ only}, not on the
angular variables. This means drag does not introduce new harmonic content
in $\bar\omega$.
\end{enumerate}


\section{Modified Mean ODE}

\subsection{Combined zonal + drag mean system}

The mean ODE of the master note (Eqs.~67--70) is extended by simply adding
the drag-averaged rates:
\begin{align}
\dot{\bar\nu}
&= F_\nu(\bar\nu, \bar g),
\label{eq:nudot_combined}\\
\dot{\bar g}
&=
\underbrace{
\sum G^{(c)}_m \cos(m\bar\omega) + \sum G^{(s)}_m \sin(m\bar\omega)
}_{\text{zonal (master note)}}
+\;
\underbrace{F_g(\bar\nu, \bar g)}_{\text{drag}},
\label{eq:gdot_combined}\\
\dot{\bar Q}
&= \text{(zonal terms, unchanged)},
\label{eq:Qdot_combined}\\
\dot{\bar\Psi}
&= \text{(zonal terms, but now $\bar\nu$ varies)},
\label{eq:Psidot_combined}\\
\dot{\bar\Omega}
&= \text{(zonal terms, but now $\bar\nu$ varies)}.
\label{eq:Omegadot_combined}
\end{align}

\subsection{What changes structurally}

\begin{enumerate}
\item \textbf{$\bar\nu$ is now a state variable}, not a constant. The mean
ODE gains a 6th variable ($\bar M$ being the 5th). This is the most
significant structural change: the zonal coefficient functions
$G^{(c)}_m(\bar\nu, \bar g, \bar Q)$, $P_0(\bar\nu, \bar g, \bar Q)$, etc.,
now have a time-varying first argument.

\item \textbf{The zonal rates acquire implicit time dependence} through
$\bar\nu(t)$, because $a = (\mu/\bar\nu^2)^{1/3}$ appears in every zonal
coefficient function. However, the functional form of the zonal coefficients
is unchanged---only the value of $\bar\nu$ at which they are evaluated
changes.

\item \textbf{The system remains autonomous}: drag depends only on the
state variables (through $a$ and $e$), not on time explicitly (assuming a
static atmospheric model).

\item \textbf{The harmonic structure is preserved}: drag adds only a
$\bar\omega$-independent term to $\dot{\bar g}$ and a new equation for
$\dot{\bar\nu}$, without introducing new Fourier harmonics.
\end{enumerate}

\subsection{Why the short-period map remains valid}

The short-period corrections of the master note are functions of the
mean state:
\begin{equation}
\bm u_1 = \bm u_1(\bar\nu, \bar g, \bar Q, \bar\Psi, \bar\Omega, \bar K).
\end{equation}
These expressions are derived from the \emph{zonal} forcing only and are
parameterized by the mean orbital elements. As long as the mean state
$(\bar\nu, \bar g, \bar Q)$ evolves slowly (which it does under drag---the
timescale for significant drag-induced changes is weeks to months for
typical LEO orbits), the short-period map evaluated at the current mean
state remains an accurate first-order correction.

The error introduced by using the zonal short-period map in the presence of
drag is of order $O(\eps_{\mathrm{drag}} \cdot \eps_{\mathrm{zonal}})$, a
cross-perturbation product. This is typically much smaller than either
perturbation individually, and is consistent with the first-order theory.

Drag-induced short-period effects (once-per-orbit drag modulation due to
the altitude variation along an eccentric orbit) are neglected. These are of
order $\eps_{\mathrm{drag}}$ in the GEqOE components and would require a
separate short-period correction. For orbits where drag is significant
(LEO), the zonal short-period corrections ($\sim J_2 \sim 10^{-3}$) are
typically much larger than the drag short-period corrections ($\sim 10^{-6}$
to $10^{-5}$ for most satellites), so the omission is justified.


\section{Atmospheric Density Considerations}

\subsection{Exponential model}

The simplest approach uses the exponential model \eqref{eq:density_exp} with
altitude-dependent scale height. This is sufficient for preliminary orbit
lifetime estimation and for the secular evolution of the mean elements over
timescales of weeks to months.

\subsection{Tabulated models}

For higher fidelity, the atmospheric density can be taken from standard
models (NRLMSISE-00, JB2008, etc.) evaluated at the mean altitude. Since
only the orbit-averaged density is needed (not the instantaneous density at
each point along the orbit), the atmospheric model evaluation is infrequent
---once per mean ODE step, not once per fast-angle step.

The modified Bessel function expansion \eqref{eq:bessel_expansion} can be
generalized to handle any density profile $\rho(r)$ by replacing the
Bessel coefficients with numerical Fourier coefficients of $\rho(r(G))$.

\subsection{Solar activity and time-dependent density}

Real atmospheric density depends on solar activity (F10.7 index, Kp index,
etc.), which introduces explicit time dependence. This breaks the autonomy
of the combined zonal + drag system. The practical solution is to treat the
solar activity indices as slowly varying parameters that are updated at each
mean ODE step, either from historical data or from prediction models.

This is analogous to how the master note's mean ODE treats the zonal
coefficients as fixed parameters: the atmospheric parameters are ``frozen''
over each integration step of the mean ODE and updated between steps.


\section{Validation Strategy}

\subsection{Reference solution}

The drag-extended mean model can be validated against full osculating
GEqOE Taylor integration (heyoka) with both zonal harmonics and drag
included. The reference integrator already supports general non-conservative
perturbations through the force model interface.

\subsection{Test cases}

\begin{enumerate}
\item \textbf{Circular LEO decay}: $a_0 = 6700$~km, $e_0 = 0.001$,
$i = 45^\circ$. Track the secular decay of $a$ and $e$ over 30 days.
Compare mean element evolution against per-orbit means extracted from the
osculating solution.

\item \textbf{Eccentric LEO with drag}: $a_0 = 7000$~km, $e_0 = 0.1$,
$i = 60^\circ$. The eccentric orbit samples a wider altitude range, testing
the Bessel expansion of the density model.

\item \textbf{Transition orbit}: $a_0 = 7500$~km, $e_0 = 0.2$,
$i = 28.5^\circ$ (GTO-like perigee). The high eccentricity concentrates drag
near perigee. This tests the regime where drag short-period effects are
largest.
\end{enumerate}

\subsection{Error metrics}

\begin{enumerate}
\item Mean element evolution: compare $\bar\nu(t)$, $\bar g(t)$ from the
drag-extended mean model against per-orbit means from the osculating
solution.

\item Osculating reconstruction: apply the zonal short-period map to the
drag-extended mean state and compare against the full osculating solution.
The reconstruction error should be dominated by the zonal $O(\eps^2)$
residual, not by the drag omission.

\item Orbit lifetime: compare the predicted time to a specified altitude
threshold against the osculating reference.
\end{enumerate}


\section{Limitations and Extensions}

\subsection{Atmospheric rotation}

The no-rotation approximation ($\bm v_{\mathrm{rel}} \approx \bm v$)
introduces an error that grows with altitude and with inclination (the
relative velocity between the satellite and the co-rotating atmosphere
depends on the satellite's velocity component in the equatorial plane).
The correction involves an additional out-of-plane drag component that
affects $Q$ and the angular variables. Including atmospheric rotation
would modify \eqref{eq:Qdot_drag}--\eqref{eq:Omegadot_drag}.

\subsection{Drag short-period corrections}

If drag short-period effects are needed (e.g., for high-fidelity osculating
reconstruction in LEO), a separate short-period correction for drag could
be derived. The drag forcing is \emph{not} a Laurent polynomial in $F$
(because of the exponential density model), so the direct residue
decomposition does not apply directly. Options include:
\begin{itemize}
\item Truncating the Bessel expansion at a finite number of terms and
treating each term as a trigonometric-polynomial forcing (amenable to the
existing machinery).
\item Numerical quadrature for the drag short-period correction (hybrid
analytical-numerical approach).
\end{itemize}

\subsection{Coupling with second-order zonal theory}

If the second-order zonal theory (companion note
\texttt{geqoe\_averaged\_second\_order}) is implemented, the drag extension
would still apply in the same way: add the drag-averaged secular rates to
the (now more accurate) mean ODE. The second-order short-period map would
provide more accurate osculating reconstruction, and the drag secular rates
would be unchanged.


\section{Implementation Roadmap}

\begin{enumerate}
\item \textbf{King-Hele drag rates}: Implement the orbit-averaged drag rates
$F_\nu(\bar\nu, \bar g)$ and $F_g(\bar\nu, \bar g)$ using the King-Hele
formulas with the exponential density model. This is a self-contained
module that takes $(a, e, C_D A/m, \rho_0, H)$ and returns the averaged
rates.

\item \textbf{Extended mean ODE}: Add $\bar\nu$ as a state variable in the
mean propagator and include the drag rates in the right-hand side. The
existing zonal coefficient functions are unchanged; they simply receive a
time-varying $\bar\nu$ argument.

\item \textbf{Validation}: Run the test cases of Section~6 against
osculating Taylor integration with drag.

\item \textbf{Atmospheric model upgrade}: Replace the exponential model
with a tabulated model (NRLMSISE-00) for higher fidelity. This affects only
the density evaluation, not the averaging procedure.

\item \textbf{Atmospheric rotation} (optional): Include the co-rotation
correction for improved accuracy at higher altitudes and low inclinations.
\end{enumerate}


\section{Summary}

The drag extension is the simplest non-conservative addition to the GEqOE
averaged theory because:
\begin{enumerate}
\item Drag affects only two of the six mean state variables ($\bar\nu$
and $\bar g$) at leading order.
\item The drag-averaged rates are known functions of $(a, e)$ that do not
introduce new harmonic content in $\bar\omega$.
\item The zonal short-period map remains valid without modification.
\item The combined zonal + drag mean system remains autonomous (for a
static atmospheric model).
\item The implementation is additive: the drag module can be developed and
validated independently of the zonal theory.
\end{enumerate}

The main limitation is the neglect of drag short-period effects, which is
justified when the drag perturbation is smaller than the zonal perturbation
in its once-per-orbit content. For typical LEO satellites, this condition
is satisfied by a factor of $\sim$100--1000.

\end{document}
